\newpage
\appendix

\section{Implementation Details}
\label{app:impl}

\subsection{Expert Artifact Corpora}
\label{app:corpora}

For each of the three domains in Section~\ref{sec:setup}, the expert
artifact $y^\ast$ is drawn from a publicly available source and used as
the anchor for retrospective reconstruction.

\begin{itemize}[leftmargin=*]
  \item \textbf{Scientific writing.}
    We use author-written related-work sections as expert artifacts
    $y^\ast$. Each instance consists of a parsed related-work block
    together with the corresponding paper abstract and cited reference
    list. The abstract is used as side context for inducing the task
    specification $\hat{x}$, while the original related-work text is
    never exposed to the agent during executable trajectory generation.

  \item \textbf{Financial analysis.}
    We use Item~7 (Management's Discussion and Analysis, MD\&A)
    sections of U.S. SEC 10-K filings from the EDGAR corpus
    \citep{loukas2021edgar}. Each MD\&A section is treated as the
    expert artifact $y^\ast$ and is indexed by company and fiscal year.
    The remaining filing content, such as business description, risk
    factors, and quantitative disclosures, is available to the agent
    only through actual filing-reader tool calls.

  \item \textbf{Legal judgment drafting.}
    We use Chinese first-instance civil and administrative judgments
    collected from publicly accessible court-record portals. Each
    judgment is normalized into plain text, with personal identifiers
    pseudonymized when necessary. Statutes and prior cases are not
    bundled with $y^\ast$ and must be retrieved at trajectory time
    through the legal search and browse tools.
\end{itemize}

Across all domains, the target artifact is used only as a retrospective
anchor for task induction, evidence recovery, and verification. It is not
available as an observation during trajectory execution.

\subsection{Procedural Diversity}
\label{app:diversity}

To prevent the reconstructed corpus from collapsing into a single
template, we randomize five axes when sampling each trajectory:
(i) user-query phrasing;
(ii) system prompt style, including variations in tone, persona, and
tool schema;
(iii) high-level plan flow pattern;
(iv) thought templates for each action type, such as \textsc{Search},
\textsc{Browse}, \textsc{Synthesize}, \textsc{Assemble},
\textsc{Reflection}, \textsc{Transition}, and \textsc{ErrorRecovery};
and (v) the surface formatting of tool-call blocks, such as JSON-style
tool calls and ReAct-style \texttt{Action/Action Input} formatting. All
random seeds are fixed within a single backbone run for reproducibility.

\subsection{SFT Mixture and Training}
\label{app:sft}

The same SFT recipe is applied to all fine-tuned variants. The
\textbf{Base} baseline denotes the original checkpoint without SFT.

\paragraph{Token-budgeted mixture.}
We treat the training corpus as a budget of $5\times10^{7}$ tokens. For
\textsc{RetroGen}, this budget is split as:
$40\%$ \textsc{Trajectory-Scientific},
$20\%$ \textsc{Trajectory-Financial},
$20\%$ \textsc{Trajectory-Legal}, and
$20\%$ \textsc{General}. The \textsc{General} pool aggregates
instruction-following, multi-turn dialogue, long-form writing, math,
and code samples. Samples longer than $32{,}768$ tokens are dropped.
The same token budget and mixture proportions are used for the
corresponding baselines, with only the $80\%$ agentic slice changed
according to the baseline definition.

\paragraph{Baseline construction.}
For \textbf{ForwardGen}, trajectories are generated from the inferred
task specification $\hat{x}$ without access to the expert artifact
$y^\ast$ as a target anchor. For \textbf{Artifact-Only}, we remove
intermediate tool calls, observations, and synthesis notes while
preserving the inferred task and final artifact. For
\textbf{Public-Baseline}, we replace the $80\%$ agentic trajectory
slice with WebGLM-QA and WebCPM-WK examples while keeping the same
$20\%$ general-data mixture. All fine-tuned variants share the same
sequence length cap and optimization configuration.

\paragraph{Optimization.}
Training is implemented with full-parameter supervised fine-tuning
using DeepSpeed ZeRO-1 and sequence packing. We use AdamW with peak
learning rate $1\times10^{-5}$, cosine scheduling, warmup ratio $5\%$,
weight decay $0.01$, gradient clipping $1.0$, gradient checkpointing,
left truncation, dropout off, and a target effective batch size of
approximately $512$K tokens per optimizer step. Each backbone is
trained for one epoch with a fixed random seed. The same optimization
configuration is used for every fine-tuned baseline.

\subsection{Evaluation Details}
\label{app:eval}

All evaluations are run with a vLLM backend.
We use \texttt{max\_length}=32{,}768 for evidence-grounded and legal
tasks and $4{,}096$ otherwise. Unless otherwise specified, the random
seed is fixed to $1234$. 

For dynamic agentic evaluation in Section~\ref{sec:astabench}, we run
AstaBench's literature-understanding suite, including
LitQA2-Validation and SQA. LitQA2 measures accuracy, precision, and
coverage on biomedical evidence-seeking questions. SQA reports
ingredient recall, answer precision, citation precision, citation
recall, and their global average. The base model and its
\textsc{RetroGen}-trained counterpart use the same backbone, tool
inventory, sampling configuration, and maximum interaction budget
(\texttt{max\_rounds}=15). Judge-based metrics are rescored using the
same local judge model for all systems.

\subsection{AstaBench Action Abstraction}
\label{app:astabench_actions}

For the transition analysis in Figure~\ref{fig:exp_trans}, we collapse
raw AstaBench tool calls into three high-level actions:
\textsc{Search}, \textsc{Read}, and \textsc{END}. \textsc{Search}
includes paper search, snippet search, citation expansion, and other
retrieval-oriented calls. \textsc{Read} includes opening paper
metadata, reading paper content, and inspecting retrieved passages.
\textsc{END} denotes the model's decision to stop tool use and produce
the final answer. We compute transition probabilities by counting each
consecutive action pair and normalizing by the total outgoing
transitions from the current action.

\subsection{Action Schema and Reusable Templates}
\label{app:schema}

For all three domains, an assistant turn may contain one or more tool
calls followed by an optional free-form note. A final-answer turn emits
no tool call. The default schema is:
\begin{lstlisting}[basicstyle=\ttfamily\scriptsize,breaklines=true]
<tool_call>{"name":"<tool>","arguments":{...}}</tool_call>
\end{lstlisting}
Tool results are returned in the next user turn:
\begin{lstlisting}[basicstyle=\ttfamily\scriptsize,breaklines=true]
<tool_response>...</tool_response>
\end{lstlisting}
We additionally sample a ReAct-style variant at fixed probability per
trajectory:
\begin{lstlisting}[basicstyle=\ttfamily\scriptsize,breaklines=true]
Action: <tool>
Action Input: {...}
\end{lstlisting}
This improves robustness to different tool-call surface forms. Thought
templates that bracket each action are drawn from a large pool of
paraphrases to avoid wording collapse during SFT.

\section{Verification Protocol}
\label{app:verification}

Verification is performed by the same backbone that generates the
candidate trace, under a rubric-guided scoring protocol; we do not
introduce a stronger external judge. Each retained trace receives four
scores in $[0,1]$---rubric satisfaction $s_{\mathrm{rub}}$, holistic
quality $s_{\mathrm{qual}}$, evidence grounding $s_{\mathrm{grd}}$, and
trace consistency $s_{\mathrm{con}}$---which are averaged with uniform
weights $\lambda=0.25$.

\paragraph{Domain-specific thresholds.}
Table~\ref{tab:score_thresholds} reports the score distribution used to
calibrate filters. Rather than applying a global cutoff, we remove the
low-score tail in each domain while retaining enough data for the
50M-token SFT mixture. The resulting thresholds are $0.52$ / $0.48$ /
$0.63$ for scientific / financial / legal traces, yielding approximately
$25$K verified trajectories.

\begin{table}[t]
\centering
\small
\setlength{\tabcolsep}{5pt}
\resizebox{\linewidth}{!}{%
\begin{tabular}{lccccc}
\toprule
\textbf{Domain} & \textbf{Mean} & \textbf{Median} & \textbf{P25} & \textbf{P75} & \textbf{Thr.} \\
\midrule
Scientific & 0.571 & 0.594 & 0.412 & 0.743 & 0.52 \\
Financial & 0.508 & 0.517 & 0.273 & 0.738 & 0.48 \\
Legal & 0.624 & 0.668 & 0.480 & 0.799 & 0.63 \\
\bottomrule
\end{tabular}%
}
\caption{Domain-wise distribution of the overall verification score and the corresponding filtering threshold. Thresholds cut the low-score tail while retaining enough traces for the 50M-token SFT mixture, yielding approximately 25K verified trajectories.}
\label{tab:score_thresholds}
\end{table}


\paragraph{Failure modes.}
Rejected traces typically exhibit missed evidence themes, weak
grounding, shallow comparison, or incomplete finals. For example, one
rejected scientific trajectory (overall score $0.30$, rubric score
$0.00$) is asked to write a related-work section on adaptive
optimization under partial observability, but produces fluent prose
about uncertainty and decision-making while omitting the required
themes. This is a common rejection pattern: a plausible-sounding
workflow that recovers the wrong evidence lineage. Accepted traces, by
contrast, search and browse in a noisy retrieval space, retain relevant
sources, discard weakly related results, and synthesize an artifact that
covers the self-induced rubric.

\section{Prompts and Anti-Leakage}
\label{app:prompts}

We include simplified prompt templates used for rubric extraction and
task induction in the scientific-writing domain. Financial and legal
domains follow the same structure with domain-specific criteria
(coverage of filings or statutes, analytical depth, and citation of
controlling authority).

\vspace{2.5pt}
\noindent\textbf{Rubric extraction.}\\[0.6ex]
The scientific-writing prompt template is:
\vspace{0.55ex}
\begin{lstlisting}
Analyze the following Related Work section and extract a structured
evaluation rubric.

## Related Work Section
{related_work}

## Instructions
Identify specific quality criteria that a reconstructed trajectory
should satisfy. Each item should be categorized as one of:
coverage | structure | synthesis | citation

Output JSON:
{
  "rubric_items": [
    {
      "id": "R1",
      "category": "coverage|structure|synthesis|citation",
      "description": "Specific requirement description"
    }
  ]
}
\end{lstlisting}
An example extracted item is:
\begin{lstlisting}
{
  "id": "R1",
  "category": "synthesis",
  "description": "Compare greedy frameworks for adaptive
optimization with Stochastic Depletion problems."
}
\end{lstlisting}

\paragraph{Task induction with anti-leakage.}
The inferred query must specify a realistic professional intent without
leaking the evidence the agent is expected to recover:
\begin{lstlisting}[basicstyle=\ttfamily\scriptsize,breaklines=true]
Based on the following Related Work section, generate a task
description for an AI agent that needs to write a similar Related
Work section using only web search.

## Related Work Section
{related_work}
## Rubric
{rubric}

## Instructions
1. Infer what paper this Related Work belongs to.
2. Specify the domain and key themes.
3. Do NOT reveal specific paper titles/authors, target citations,
   exact conclusions, or evidence that the agent must recover
   via search.
4. Provide enough high-level guidance for search.
\end{lstlisting}
A representative inferred task is: \emph{write a related-work section on
trust-aware recommender systems and matrix factorization, covering
fundamental factorization algorithms, social-network structure, and
trust propagation.} The prompt names the research intent and topical
scope, but still requires the agent to search, inspect, select, and
synthesize evidence. During trajectory execution the original artifact
$y^{*}$ is used only offline as a retrospective anchor.
